\bookchapter{Kapitel 31: Einstieg in Kurzvideo-Drehs und Content-Erstellung}

\chapterintro{2026 sind Kurzvideos kein “Spielzeug der Jungen” mehr – sie sind die grundlegende Art, wie alle Menschen Informationen erhalten und sich ausdrücken. Ob Sie das Leben festhalten, Wissen teilen oder einen Nebenverdienst aufbauen wollen – das Handy ist Ihr bestes Werkzeug. Keine Kamera, kein Premiere nötig. Ein Handy genügt. Dieses Kapitel hilft Ihnen beim Start von Null.}

\booksection{I. Vor dem Dreh überlegen: Für wen drehen Sie?}

Der erste Fehler vieler: “Handy in die Hand nehmen und loslegen” – ohne zu wissen, für wen. Bevor Sie drehen, beantworten Sie drei Fragen:

\begin{enumerate}
  \item \textbf{Wer ist mein Publikum?}—— Für Kinder, für Gleichaltrige oder für ältere Menschen – der Stil ist völlig unterschiedlich.
  \item \textbf{Welchen Wert biete ich?}—— Zum Lachen bringen (Unterhaltung), lernen lassen (Wissen), Appetit machen (Essen), berühren (Emotion). Mindestens eines.
  \item \textbf{Auf welcher Plattform kann ich regelmäßig posten?}—— Douyin/Kuaishou (Unterhaltung), Xiaohongshu (Lifestyle), Bilibili (tiefe Inhalte), Video-Kanal (sozialer Kreis). Wählen Sie eine Hauptplattform und vertiefen Sie sich – nicht sofort auf allen veröffentlichen.
\end{enumerate}

\booksection{II. Grundlegende Handy-Einstellungen —— 5 Pflichtaktionen vor dem Dreh}

\begin{enumerate}
  \item \textbf{Auflösung auf 1080P, Bildrate auf 30 fps}: Das ist der Standard für Kurzvideoplattformen. 4K ist zu groß und lädt langsam, 60 fps bringt auf Social Media keinen Vorteil, macht die Datei aber größer.
  \item \textbf{Objektiv sauberwischen}: Die erste und wichtigste Regel. Vor dem Dreh das Objektiv mit dem Ärmel abwischen – das Bild wird sofort deutlich schärfer.
  \item \textbf{Fokus und Belichtung sperren}: Im Sucher 2 Sekunden lang auf das Hauptmotiv drücken, bis "AE/AF-Sperre" erscheint. So wird das Bild beim Bewegen nicht ständig hell und dunkel.
  \item \textbf{Quer- oder Hochformat?}—— Für Douyin/Kuaishou/Video-Kanal/Xiaohongshu\textbf{Hochformat}(9:16). Für Bilibili/YouTube\textbf{Querformat}(16:9). Erst Plattform überlegen, dann Format wählen.
  \item \textbf{Gitterlinien einschalten}: Einstellungen > Kamera > Gitter. Die Grundlage der Drittel-Regel, hilft, schnell den Bildschwerpunkt zu finden.
\end{enumerate}

\booksection{III. Fünf Kamera-Bewegungen —— das Bild in Bewegung versetzen}

Das größte Problem von Anfängern ist zu “ruhiges” Bild – die ganze Zeit stillstehend gefilmt. Lernen Sie diese fünf Grundbewegungen, und das Bild wirkt sofort professioneller:

\begin{itemize}
  \item \textbf{Push (Heranfahren)}: Handy nähert sich langsam dem Motiv. Erzeugt Spannung oder Nähe.
  \item \textbf{Pull (Wegfahren)}: Handy entfernt sich langsam vom Motiv. Vom Detail zur Totale – zeigt die Umgebung.
  \item \textbf{Pan (Schwenken)}: Handy bleibt an Ort, dreht sich links/rechts. Gut für Landschaften und Räume.
  \item \textbf{Track (Fahrt)}: Handy bewegt sich parallel, folgt dem Motiv. Braucht ruhige Hand, am besten mit Gimbal.
  \item \textbf{Follow (Verfolgen)}: Handy geht hinter oder vor dem Motiv her. Die häufigste Bewegung bei Vlogs.
\end{itemize}

\textbf{Ruhe ist das Wichtigste}: Wackelige Bilder wirken unprofessionell. Beide Hände sind ruhiger als eine; beim Gehen Knie leicht beugen, “Ninja-Schritt” zur Reduzierung von Stößen; bei Budget einen Gimbal kaufen (DJI OM-Serie ca. 500–800 Yuan) – qualitative Verbesserung.

\booksection{IV. Videoschnitt für Einsteiger —— CapCut reicht}

CapCut (die chinesische Version: JianYing) ist aktuell das beste Schneideprogramm für Anfänger – kostenlos, funktionsreich, viele Vorlagen. Der Grund-Workflow:

\begin{enumerate}
  \item \textbf{Material importieren}: CapCut öffnen > Neues Projekt > Videoclips auswählen und importieren.
  \item \textbf{Schneiden}: Zeitleiste verschieben und an Anfang und Ende jeder Szene überflüssiges Material abschneiden. Prinzip: Nur informative Teile behalten – lieber kürzer als zu lang.
  \item \textbf{Übergänge hinzufügen}: Zwischen zwei Clips "Ubergang" antippen und einen einfachen Effekt wählen. Anfängern wird "Blend" oder "Schwarzblende" empfohlen – keine ausgefallenen 3D-Übergänge.
  \item \textbf{Untertitel hinzufügen}: CapCut hat die Funktion "Untertitel erkennen" – automatisch generierte Untertitel mit einem Klick. Genauigkeit ca. 90 \%, einmal manuell korrigieren genügt.
  \item \textbf{Musik hinzufügen}: CapCut hat eine große Bibliothek lizenzfreier Musik. Wählen Sie ein BGM, das zum Rhythmus und Inhalt passt, und setzen Sie die Lautstärke auf ca. 1/3 des Originaltons.
  \item \textbf{Exportieren}: In 1080P exportieren, 30 fps, Bitrate "empfohlen" wählen.
\end{enumerate}

\booksection{V. Die Bedeutung des Thumbnails —— 70 \% entscheiden nur anhand des Covers}

Das Thumbnail eines Kurzvideos ist die “Schaufenster” – es entscheidet, ob jemand klickt. Wichtige Punkte:

\begin{itemize}
  \item \textbf{Das Thumbnail braucht einen Satz}: Schreiben Sie den Kerninhalt in großen Buchstaben, z. B. "In 3 Minuten Handy-Fotos bearbeiten". Große Schrift, helle Farben, Kontrast zum Hintergrund.
  \item \textbf{Thumbnail und Inhalt müssen übereinstimmen}: Keine "Clickbait" – wenn das Thumbnail etwas anderes verspricht als der Inhalt, scrollen die Zuschauer sofort weg und der Algorithmus stuft Sie ab.
  \item \textbf{Lebendige Gesichtsausdrücke}: Wenn das Thumbnail eine Person zeigt, wählen Sie einen lebendigen oder übertriebenen Moment. Gesichter ziehen mehr Klicks an als Text.
  \item \textbf{Thumbnail mit CapCut}: In CapCut das beste Bild wählen, Text überlagern und als Bild exportieren.
\end{itemize}

\booksection{VI. Drei Prinzipien der Content-Erstellung}

\begin{itemize}
  \item \textbf{Die ersten 3 Sekunden entscheiden}: Auf Kurzvideoplattformen entscheiden Zuschauer in 3 Sekunden. Die ersten 3 Sekunden brauchen einen "Hook" – eine Frage, einen Cliffhanger, ein überraschendes Bild.
  \item \textbf{Ein Video, ein Thema}: Der häufigste Anfängerfehler: zu viel Inhalt in ein Video. Ein Video löst ein Problem. Gründlich ist wichtiger als umfassend.
  \item \textbf{Durchhaltevermögen ist wichtiger als Talent}: Der Empfehlungsalgorithmus braucht "Account-Pflege". Die ersten 20 Videos haben fast keine Aufrufe – das ist normal. Wöchentlich 2–3 Videos hochladen, damit der Algorithmus Ihre Inhaltsrichtung erkennt; der Traffic kommt langsam.
\end{itemize}

\begin{tcbpractice}{Praxis: Ein 30-Sekunden-Kurzvideo drehen}
\begin{enumerate}
  \item Wählen Sie ein kleines Thema, das Sie gut kennen (z. B. "Wie man die Handy-Schrift vergrößert").
  \item 3–5 Clips im Hochformat drehen, jeder 5–10 Sekunden.
  \item Mit CapCut zusammenfügen, schneiden, Untertitel und BGM hinzufügen.
  \item Ein Thumbnail mit großer Schrift erstellen.
  \item Nach dem Export Freunden zeigen, Feedback einholen, dann veröffentlichen.
\end{enumerate}

\end{tcbpractice}

\begin{tcbsummary}{Kapitelzusammenfassung}
\begin{itemize}
  \item Vor dem Dreh überlegen: Wer ist das Publikum, welchen Wert bieten Sie, auf welcher Plattform
  \item 5 Dinge vor dem Dreh: 1080P/30 fps, Objektiv wischen, Fokus/Belichtung sperren, Format wählen, Gitter an
  \item Fünf Bewegungen: Push, Pull, Pan, Track, Follow – Ruhe zuerst
  \item Schnitt mit CapCut: Schneiden Übergang Untertitel BGM Export, sechs Schritte
  \item Erste 3 Sekunden entscheiden, ein Video ein Thema, Durchhaltevermögen schlägt Talent
\end{itemize}

\end{tcbsummary}
